\section{Related Work}
\label{sec:related_work}

\textbf{Vulnerability/Bug Detection and Prevention}. Various methods have been proposed to determine if a code component (component, file, function/method, or statement/line) is vulnerable. The rule-based techniques apply static analyzers and leverage seen vulnerability patterns, such as FlawFinder~\cite{FlawFinder} or Coverity~\cite{Coverity}. Recently, several deep-learning based approaches have been introduced~\cite{survey_papers, poster, vuldeeppeaker, vulsniper, devign, are_we_there,mvd,issta_22}. 
VulDeePecker~\cite{vuldeeppeaker} and SySeVR~\cite{sysevr} introduce tools to detect \textit{slice-level} vulnerabilities, which are more fine-grained.
IVDetect~\cite{ivdetect}, which is a graph-based neural network model, is proposed to detect vulnerabilities at the function level and use a model interpreter to identify vulnerable statements in the detected suspicious functions. LineVul~\cite{linevul} and LineVD~\cite{linevd} apply CodeBERT in their own way and have been shown that they are more effective than IVDetect in detecting vulnerable functions and lines/statements. VelVet~\cite{velvet} builds graph-based models to detect vulnerable statements. Our work could complement well with the existing automated vulnerability detection approaches. Particularly, the type identification step could be applied after developers use a vulnerability detection method to quickly interpret and fix the vulnerable functions.
%
Our work could also be applied to determine the vulnerability types of a vulnerable code component before applying one or more vulnerability detection approaches specialized for certain vulnerability types, such as the approach to memory-related vulnerabilities~\cite{mvd}.
%


\hl{Our work might related VUDENC}~\cite{VUDENC} \hl{Wartschinski \textit{et al.} and $\mu$VulDeePecker}~\cite{mu_vuldeeppeker} \hl{by Zou \textit{et al.} which focus on vulnerability detection and can implicitly indicate the type of vulnerability. In other words, given a piece of code, these techniques determine if the code is benign or belongs to one/some vulnerability types. VUDENC}~\cite{VUDENC} \hl{uses separate LSTM binary classification models to determine if a piece of code is neutral/benign or belongs to some vulnerability types. 
%
%
In this paper, we focus on the specific task of vulnerability type identification integrated after vulnerability detection to provide developers with the type of vulnerabilities effectively identified by the existing vulnerability detection techniques}~\cite{vuldeeppeaker,sysevr,ivdetect,linevul,linevd} \hl{which determine if the given code is benign or vulnerable. }
%
%
Meanwhile, the goal of $\mu$VulDeePecker~\cite{mu_vuldeeppeker} is to determine if a function is clean (not vulnerable) or of exactly one vulnerability. However, our work differs from $\mu$VulDeePecker in three fundamental aspects. First, $\mu$VulDeePecker is a vulnerability detection approach which is designed to decide if a function is vulnerable. In other words, the input of that approach is a function that has not known whether it is vulnerable or not. Meanwhile, our work is designed to apply as a step after detecting vulnerable functions. Therefore, the input function of our work is assumed to be vulnerable. Moreover, $\mu$VulDeePecker assumes that a vulnerable function has only one type. Thus, the multi-class classification is considered in $\mu$VulDeePecker. Meanwhile, a function could belong to multiple types. Hence, in our work, the VTI problem is modeled as the multi-label classification task. Finally, $\mu$VulDeePecker is designed for the function call vulnerabilities, while our work has no limit to the kind of vulnerabilities.

% \textbf{Code Representation Learning}. 

\textbf{Learning-based approaches for SE tasks}. Several studies have been proposed for specific SE tasks, including code suggestion/completion~\cite{icse20, naturalness, autosc, apsec21,arist}, program synthesis~\cite{gvero2015synthesizing}, pull request description generation~\cite{hu2018deep,liu2019automatic}, code summarization~\cite{iyer2016summarizing,mastropaolo2021studying,wan2018improving}, code clones~\cite{li2017cclearner}, fuzz testing\cite{godefroid2017learn}, code-text translation~\cite{ase22}, and program repair~\cite{jiang2021cure,ding2020patching}.
%
Recently, several learning techniques have been proposed to learn representing source code for specific SE applications~\cite{oppsla19, code2vec} or general SE tasks~\cite{codebert,cubert,infercode}.